\section{A shared leading residue pool fails even with node biases}

This checkpoint concerns one proposed inductive invariant for the paired Haar
decoder.  It is not a lower bound for the original xorformer model, nor a
counterexample to the full paired decoder.

Let $K_{ij}=(1-p_i-p_j)^{-2}$.  A shared leading residue vector $v$ for two
child words $\sigma,\tau\in\{-1,1\}^q$ would require real biases $b,c$ with
\[
 \sigma_i((K\operatorname{diag}(\sigma)v)_i+b)>0,
 \qquad
 \tau_i((K\operatorname{diag}(\tau)v)_i+c)>0
 \quad(0\leq i<q).
\]
The entries of $v$ are unrestricted real numbers.  Separate biases therefore
do not rescue this particular shared-pool invariant.

\paragraph{Kernel-checked finite counterexample.}
Take $q=5$, $p_i=i/20$,
$\sigma=(1,-1,-1,1,-1)$ and $\tau=(1,-1,1,-1,1)$.
Eliminate the bias by subtracting a negative query from a positive query in
each word.  The six rows with indices
\[
 (\sigma,0,2), (\sigma,3,1), (\sigma,3,4),
 (\tau,0,3), (\tau,2,1), (\tau,4,1)
\]
have a strictly positive integer dependence.  The coefficients are recorded
as \texttt{jointFamilyPoolWeights} in
\texttt{JointFamilyBiasedCauchyPool.lean}; the exact identity is
\texttt{jointFamily\_biased\_pool\_identity}.  Consequently their six
evaluations cannot all be positive.  The declaration
\texttt{jointFamily\_no\_biased\_cauchy\_pool} proves the impossibility
with arbitrary real $v,b,c$.

\paragraph{The original compact pole hierarchy also fails.}
For $q=8$, let
$p_y=y_0\varepsilon+y_1\varepsilon^2+y_2\varepsilon^4$ and use
\[
 \sigma=(1,1,1,1,-1,1,-1,-1),\qquad
 \tau=(1,-1,-1,-1,1,1,1,1).
\]
Write $R_{w,a,b}=(K_{a,*}-K_{b,*})\operatorname{diag}(w)$.
The nine rows
\[
 R_{\sigma,0,6}, R_{\sigma,0,7}, R_{\sigma,1,4},
 R_{\sigma,2,4}, R_{\sigma,2,6}, R_{\sigma,3,4},
 R_{\tau,0,3}, R_{\tau,5,1}, R_{\tau,5,2}
\]
have a positive dependence for every sufficiently small
$\varepsilon>0$.  Define its last weight to be $1$ and solve the first
eight transposed columns for the remaining weights over
$\mathbb Q(\varepsilon)$.  Exact Gaussian elimination gives the leading
terms
\[
 \left(\frac{135}{7},\frac{8}{3\varepsilon},1,
 \frac{128}{7},\frac{8}{3\varepsilon},\frac{8}{3\varepsilon},
 1,\frac{8}{3\varepsilon},1\right).
\]
All are positive.  The eight pivot leading terms are nonzero, so the
rational functions are defined on a punctured neighborhood of zero.
Their defining equations give an exact zero linear combination of the
nine rows.  Each row would have strictly positive evaluation under the
claimed inequalities, a contradiction.

The file \texttt{JointFamilyBiasedPoolGerm.py} independently reconstructs
the rational-function linear system and verifies its pivot and weight
leading terms using exact truncated formal series, with 332 coefficients
of final precision.  This is an exact finite computation supporting the
written rational-function proof, not a kernel-checked Lean theorem.

\paragraph{What remains possible.}
The full paired Haar system has separate residues on the two poles in
each split pair.  Their sums can diverge as the split gap tends to zero,
creating additional higher-order terms.  The obstruction above assumes
one common leading residue difference; it does not constrain those
extra terms.  It also leaves label-dependent pole locations, different
query times, and different simultaneous decoding constructions open.
